% ============================================================================
% sections/mechanisms.tex -- Section 4: Mechanisms (compressed)
% ============================================================================

\section{Mechanisms}
\label{sec:mechanisms}

All mechanisms implement one skeleton. The service polls a batch $B$ of source records and
performs up to three durable duties: \emph{D-effect} appends the batch's postings to the
ledger sink; \emph{D-progress} makes the batch's consumption progress durable;
\emph{D-publish} appends derived entry records to the publication stream. A mechanism is
defined by (i) whether D-effect and D-progress are one atomic commit, (ii) the ordering of
the two when they are not, (iii) whether the effect is idempotent under the business key, and
(iv) how D-publish is made recoverable. The six mechanisms and three combinations below are
the points of that design space that practitioners deploy; the naive and Kafka-internal
variants are controls. The experiment measures both the correctness consequence and the cost,
which is what RQ2 asks for.

Throughout, $B$ denotes a poll batch, $\Pi(e)$ the expected posting pair of event $e$,
$\mathit{insert\_plain}$ a plain \texttt{INSERT}, and $\mathit{insert\_idem}$ an insert that is
a no-op when $(\mathit{event\_id},\mathit{leg\_no})$ already exists
\cite{pg-onconflict}. Every mechanism uses one database transaction per poll batch; the
per-record-autocommit variant is analysed as crash point $C_{12}$ but not deployed.

\subsection{The mechanisms}
\label{sec:mechanisms:defs}

\paragraph{M0 --- naive at-least-once (negative control).}
M0 commits the postings with a plain insert, publishes on a best-effort basis, and only then
commits offsets; it is the mechanism a team writes first: correct in the absence of faults,
and the reference point against which the others must improve. Its ledgers have no
de-duplication constraint, so any redelivery between the effect commit and the offset commit
duplicates postings.

\paragraph{M1 --- Kafka exactly-once, two sink variants.}
M1 (Algorithm~\ref{alg:m1}) performs the Kafka-internal appends inside a transaction
\cite{kip98,kip447,kafka-design}. In \emph{M1-kafka} the ledger \emph{is} the publication
stream: entries appended inside the transaction are the ledger, so durable effects and offsets
are one atomic Kafka commit and no external store is involved --- the positive control of the
paper. In \emph{M1x} the same transaction wraps the stream appends, but the ledger rows are
written by a separate, non-atomic database commit between the produce step and the transaction
commit. This is the deployment the paper is about: practitioners keep Kafka's exactly-once for the
internal stream and assume it covers the external write. It is not a straw man --- the
transaction is the documented way to obtain internal exactly-once, and the external write is
typically added as an ordinary database call in the same method: precisely the non-atomic
ordering of Algorithm~\ref{alg:m1}.

\begin{algorithm}[t]
\footnotesize
\begin{algorithmic}[1]
\State \textbf{on assignment} $p$: seek to committed offset
\Loop
  \State $B \gets \Call{Poll}{ }$
  \State \Call{TxnBegin}{ }
  \ForAll{$r \in B$}
    \State \Call{TxnSend}{stream, entry($r$)} \Comment{inside $T$}
  \EndFor
  \If{DB sink (M1x)}
    \State \Call{DbBegin}{}; insert postings of $B$; \Call{DbCommit}{} \Comment{outside $T$}
  \EndIf
  \State \Call{SendOffsetsToTxn}{$B$} \Comment{inside $T$}
  \State \Call{TxnCommit}{ } \Comment{stream and offsets atomic}
\EndLoop
\end{algorithmic}
\caption{M1 family: Kafka transaction (M1-kafka and M1x)}
\label{alg:m1}
\end{algorithm}

\paragraph{M2 --- transactional outbox with inbox de-duplication.}
M2 (Algorithm~\ref{alg:m2}) makes the database transaction the atomic unit. Inside it, the
mechanism claims each event in an \emph{inbox} table whose primary key is
$\mathit{event\_id}$, inserts the postings only for successfully claimed events, and writes an
\emph{outbox} row describing the publication intent. The database commit is therefore
simultaneously the effect, the de-duplication record and the publication intent. Offsets are
committed afterwards; redelivery finds the claim and becomes a no-op. A relay thread reads
unmarked outbox rows and publishes them, marking each row after the broker acknowledgement.
This is the outbox pattern as documented for microservice integration
\cite{debezium-outbox,laigner2021,micro2026}; it buys recoverable publication at the cost of
two extra durable writes per event and a background thread. (Sagas
\cite{sagas1987,saga2019} compensate rather than make commits atomic and are outside the
pipeline model studied here.)

\begin{algorithm}[t]
\footnotesize
\begin{algorithmic}[1]
\State \textbf{consumer loop}
\Loop
  \State $B \gets \Call{Poll}{ }$
  \State \Call{DbBegin}{ }
  \ForAll{$r \in B$}
    \State \Call{ClaimInbox}{r} $\rightarrow \mathit{claimed}$ \Comment{claim if absent}
    \If{$\mathit{claimed}$}
      \State \Call{insert\_plain}{$\Pi(\mathrm{parse}(r))$}
      \State \Call{AppendOutbox}{$r$}
    \EndIf
  \EndFor
  \State \Call{DbCommit}{ } \Comment{effect, inbox and outbox atomic}
  \State \Call{CommitOffsets}{$B$}
\EndLoop
\Statex
\State \textbf{relay thread}
\Loop
  \State $R \gets \Call{UnpublishedOutbox}{ }$
  \ForAll{$o \in R$}
    \State \Call{Publish}{$o$}; on ack \Call{MarkPublished}{$o$}
  \EndFor
\EndLoop
\end{algorithmic}
\caption{M2: transactional outbox with inbox claim and relay}
\label{alg:m2}
\end{algorithm}

\paragraph{M3 --- idempotency key with a unique index.}
M3 adds a unique index on the pair (event identifier, leg number) and inserts every leg with
a conflict-ignoring insert (\texttt{ON CONFLICT DO NOTHING}); a redelivery is absorbed by the
index, so the effect becomes idempotent without any additional table, and offsets are committed
after the database commit. The mechanism is the cheapest de-duplication known in practice ---
idempotency keys are the standard recipe for safe retries in public
APIs~\cite{stripe-idempotency} --- but, as Section~\ref{sec:formal} shows, it
protects against duplicates only: a posting never written is not restored by a unique index.

\paragraph{M4 --- offsets in the same transaction as the postings.}
M4 (Algorithm~\ref{alg:m4}) removes Kafka from the recovery path for correctness: the database
transaction contains both the postings and the progress row $\mathit{next\_offset}$ per
partition. On assignment --- including after every rebalance \cite{kip429} --- the consumer
seeks to the stored offset, not the Kafka-committed one. Kafka offsets may still be committed
for observability, but they are never consulted for recovery. No de-duplication index is
needed for the crash points that keep progress behind the effect, because a durable effect
implies a durable progress row; the open case is a second writer of the same partition
($C_7$). Seeking on assignment makes M4 recoverable: the group protocol can move a partition at any
moment \cite{kip429}, and a mechanism whose recovery point is a Kafka offset cannot know
whether the transaction that wrote the last batch also advanced it. M4 makes the database
transaction the only recovery point, and the monotone guard prevents a late writer from moving
the offset backwards.

\begin{algorithm}[t]
\footnotesize
\begin{algorithmic}[1]
\State \textbf{on assignment} $p$: $o \gets$ \Call{StoredOffset}{p}; seek $p$ to $o$
\Loop
  \State $B \gets \Call{Poll}{ }$
  \State \Call{DbBegin}{ }
  \ForAll{$r \in B$}
    \State \Call{insert\_plain}{$\Pi(\mathrm{parse}(r))$}
  \EndFor
  \ForAll{partitions $p$ in $B$}
    \State \Call{AdvanceProgress}{$p,\ \mathrm{end}(B,p)$} \Comment{monotone guard}
  \EndFor
  \State \Call{DbCommit}{ } \Comment{effect and progress atomic}
  \State \Call{Publish}{$B$}
\EndLoop
\end{algorithmic}
\caption{M4: offsets stored in the posting transaction}
\label{alg:m4}
\end{algorithm}

\paragraph{Combinations.}
Three combinations complete the design space for RQ3. \emph{M2+M3} replaces M2's inbox claim
by the posting unique index: the outbox is still written, but the de-duplication record is the
posting itself. \emph{M4+M3} adds the unique index to M4, closing the concurrent-writer
exposure. \emph{M1x+M3} adds the unique index to M1x's external write, keeping Kafka-internal
exactly-once for the publication stream. Other combinations are dominated in the model: M2+M4
closes no exposure that M4 does not already close, M0+M3 is M3, and three-way combinations add
cost without closing an additional exposure class.

\subsection{Distinctness, ordering, and failure handling}
\label{sec:mechanisms:compare}

The mechanisms are genuinely distinct: no two configurations share the same combination of
atomicity, idempotence, publication path and recovery source of truth. The comparison is the
bridge to Section~\ref{sec:formal}: each configuration closes a different subset of the two
windows of Eqs.~\eqref{eq:w1}--\eqref{eq:w2}.




The design space is summarised by four properties: \emph{atomicity} (effect and progress commit
together), \emph{idempotence} (the effect is deduplicated on the pair event identifier/leg
number), \emph{publication path} (inline, or inside the Kafka transaction, or through an outbox
relay) and \emph{recovery source} (Kafka offsets, the Kafka transaction, or a database progress
row). M0 and M1x are neither atomic nor idempotent; M1-kafka is atomic at the Kafka site with an
implicit transaction-level dedup; M2 is atomic at the database with an inbox primary key and an
outbox relay; M3 is idempotent through a posting unique index; M4 is atomic at the database
through a progress row. The combinations add the missing property to their base mechanism:
M2+M3 replaces the inbox key by the posting index, M4+M3 adds the index against the concurrent
writer, and M1x+M3 adds it to the external write while keeping Kafka-internal exactly-once for
the publication stream.

\paragraph{Commit ordering.}
For mechanisms that do not commit both duties at one site, the program order of
$\mathrm{Commit}_{S}$ (effect) and $\mathrm{Commit}_{K}$ (progress) selects which window of
Lemma~\ref{lem:windows} stays open. All mechanisms are deployed in \emph{strict} order ---
effect before progress --- so $W_2$ is out of reach and $W_1$ is the exposure that
de-duplication must close. Because H3 concerns the loss window, every mechanism is additionally
instrumented with an \emph{early} variant that commits progress before the effect; early is a
probe, not a candidate design, and exists to make $W_2$ observable under a deterministic crash
point ($C_5$, class $K_4$). The probe uses the same code path with a single ordering flag, so it
cannot be dismissed as a code mutation; the bounded model of Section~\ref{sec:formal} reproduces
both orderings.

\paragraph{Failure handling.}
The retry policies form part of each mechanism's definition: they determine whether an unknown
commit outcome is re-established by re-execution or by checking the de-duplication state. After a
definite failure every mechanism retries (M0/M1x up to three times with backoff, then replay);
after an unknown outcome only a mechanism that can re-establish the outcome from its
de-duplication state or stored progress may retry safely --- M2 and M3 do so through the inbox
and the unique index, M4 re-reads stored progress and skips if it advanced, while M0/M1x must
replay without a write retry. Kafka-offset or transaction-commit failures are retried by
M0/M1x/M2/M3 and are best effort for M4, which never uses Kafka offsets for recovery; a fenced
M1 transaction aborts and restarts the instance. Fixing the policy as part of the mechanism
definition prevents a comparison in which one mechanism benefits from a more optimistic retry
rule than another.

\paragraph{Fairness contract.}
All mechanisms share the same client, broker, database and JVM settings; the configuration is
hashed per run and an audit rejects any run whose shared part differs (Section~\ref{sec:method}).
Only mechanism-intrinsic components may differ: M1's transactional producer, M2's inbox/outbox
tables and relay, the de-duplicating index, and M4's progress table. No mechanism receives
per-mechanism tuning; a configuration change found during the pilot is applied wherever it
applies and recorded as a deviation. This contract makes the performance comparison a comparison
of mechanisms, not of tuning effort.
